% sec-05-discussion.tex - Discussion (full)
\section{Discussion}

\subsection{Probable Benefit Exceeds Probable Risk}

The central claim is that, for a Phase 1 PDAC population, the probable benefit of
repository based LLM document generation exceeds its probable risk. The benefits are
concrete and on the critical path: faster paperwork, faster treatment, a compressed
administrative and preparation bucket, and a larger number of figures available for
human oversight than a manual process would produce. The risks are the known LLM
failure modes, and each is bounded by a specific control rather than left to trust.
An occasional generation or server error is caught by human-in-the-loop review at
every commit; a formatting artifact is caught because each figure is grounded in a
Mermaid source and each file is visible on GitHub; an unverified claim is caught by
the name-matching and proofreading passes. The residual risk after these controls is
low, while the benefit is realized on every iteration. For patients whose disease
will not pause, the alternative (continuing to author large documents by slower
sequential means) carries its own, often larger, cost in delayed treatment.

\subsection{Mechanism: Compressing Administrative and Preparation Time}

The mechanism is specific. A clinical trial timeline divides into clinical and
operational time, administrative and preparation time, and fixed regulatory review
time \cite{srcDocTypesGem}. Writing faster does not make a tumor respond faster and
does not compel a regulator to review faster; it compresses only the middle bucket.
That bucket, however, contains the document white space between phases, which is
often where months are lost. By generating the six highest-value documents
prospectively and in synchronized form (so that protocol, consent, database, and
site documents move together rather than serially), the method removes precisely the
idle authoring time that separates the end of one step from the start of the next.

\subsection{Prior Developments as Evidence for This Project}

This paper did not arrive without precedent. The author's 2024 to 2026 works
established, in stages, that LLMs can be trusted for oncology trial artifacts: from
drug-cost reasoning and clinical decision support, through glioblastoma meta-analysis
and lung adenocarcinoma reviews \cite{gbmmeta10yr, lungadeno}, to PDAC digital-twin
proposals and large in silico trials \cite{kawchak_2025_15735068,
kawchak_2025_16415815, kawchak_2025_17001137}, to compliance-aware efficiency and
verified comparative code generation \cite{e2eoncoeff, codegen16}. The National
Platform for Physical AI Oncology Trials provided the governing framework
\cite{natplatform}, and the Phase 1 \cite{phase1protocol} and Phase 2
\cite{phase2protocol} protocols, together with the law and bill repositories
\cite{lawnarrative, adoptionguide, lawenactment, hr9510v5, bill01v4}, demonstrated
the single-prompt multi-file build on real regulatory and legislative documents.
Applying the same single-prompt approach to the document-generation problem itself is
the natural next step.

\subsection{Why Trust the Method: Verification, Not Faith}

The method is designed around verification before generation \cite{hr9510v5}: a
document is trusted because it is grounded, name-matched, monitorable, and quality-
gated, not because the model is assumed correct. This matters as AI outputs grow more
complex and voluminous. The adoption guide frames the underlying question politely
but pointedly: as AI-directed artifacts become larger and more intricate, why would
human-only review remain sufficient, rather than a discipline that makes every claim
traceable to a source and every figure traceable to a Mermaid file? Figure 21 quality gates are operational answer.

\subsection{Relationship to the Real-World Daraxonrasib Trial}

The documents accelerated here are exactly those the real daraxonrasib PDAC program
requires. Figure 23 traces the document thread from a KRAS-mutated PDAC patient and
the daraxonrasib (RMC-6236) plus LLM-directed robotic Whipple intervention, through
the Phase 1 combined IND/IDE protocol \cite{phase1protocol} and the 3+3
dose-escalation cohort-review packages that establish the recommended Phase 2 dose,
to the amendments, the Phase 1-to-2 transition documents, and the randomized Phase 2
protocol \cite{phase2protocol, kawchak_2026_20196639}. Figure 24 closes the argument:
faster authoring shortens the white space between these documents, which advances
treatment and, for this population, extends progression-free and overall survival.

\begin{mermaidfig}
\node[mmin] (pt) at (0.2,0) {KRAS-mutated PDAC\\daraxonrasib + robotic Whipple};
\node[mmproc] (p1) at (4.7,0) {Phase 1 IND/IDE protocol};
\node[mmproc] (p2) at (9.2,0) {3+3 cohort review\\establish RP2D};
\node[mmproc] (p3) at (9.2,-1.8) {Amendments +\\synchronized consent};
\node[mmproc] (p4) at (4.7,-1.8) {Phase 1-to-2 CSR\\and EOP2 briefing};
\node[mmgoal] (p5) at (0.2,-1.8) {Phase 2 randomized\\protocol};
\node[mmaccent] (acc) at (4.7,-3.7) {each document built by the single-prompt workflow};
\draw[mmedge] (pt)--(p1); \draw[mmedge] (p1)--(p2);
\draw[mmedge] (p2)--(p3); \draw[mmedge] (p3)--(p4); \draw[mmedge] (p4)--(p5);
\draw[mmedged] (acc.north) to (p4.south);
\end{mermaidfig}
\vspace{-0.7cm}
\figcaption{Figure 23. The real-world daraxonrasib PDAC document thread, from patient
and\\intervention through the Phase 1 protocol, cohort-review packages, amendments,
and\\transition documents to Phase 2 randomized protocol, each built by
single-prompt workflow.}

\begin{mermaidfig}
\node[mmaccent] (a) at (0,-0.2) {Faster validated\\authoring};
\node[mmproc] (b) at (4.4,-0.2) {Earlier IND/IRB,\\site activation};
\node[mmproc] (c) at (8.8,-0.2) {Earlier first dose,\\faster cohorts};
\node[mmproc] (d) at (8.8,-2.0) {Shorter white space\\between phases};
\node[mmgoal] (e) at (4.4,-2.0) {Earlier access to\\effective therapy};
\node[mmgoal] (f) at (0,-2.0) {Extended PFS / OS};
\node[mmctx] (g) at (0,-3.8) {does not change tumor biology or fixed review clocks};
\draw[mmedge] (a)--(b); \draw[mmedge] (b)--(c); \draw[mmedge] (c)--(d);
\draw[mmedge] (d)--(e); \draw[mmedge] (e)--(f);
\draw[mmedged] (g.north) to (f.south);
\end{mermaidfig}
\vspace{-0.7cm}
\figcaption{Figure 24. The patient time-saved cascade. Faster authoring advances
each step from\\submission to first dose, shortening the white space between phases
and, for this population,\\extending progression-free, overall survival, without
changing tumor biology or fixed review clocks.}
